\subsection{DIRECT LIMITS}

Let I be a (right) directed preordered set and let \((\mathbf{E}_{\alpha})_{\alpha\in\mathbf{I}}\) be a family of sets indexed by I. For each pair \((\alpha,\beta)\) of elements of I such that \(\alpha\leqslant\beta\) , let \(f_{\beta\alpha}\) be a mapping of \(E_{\alpha}\) into \(E_{\beta}\) . Suppose that the \(f_{\beta\alpha}\) satisfy the following conditions:

(LI \(_{I}\) ) The relations \(\alpha \leqslant \beta \leqslant \gamma\) imply \(f_{\gamma\alpha} = f_{\gamma\beta} \circ f_{\beta\alpha}\) .

(LI \(_{II}\) ) For each \(\alpha \in I\) , \(f_{\alpha\alpha}\) is the identity mapping of E \(_{\alpha}\) .

Let G be the set which is the sum of the family of sets \((\mathbf{E}_{\alpha})_{\alpha\in\mathbf{I}}\) (Chapter II, §4, no. 8); by abuse of language, we shall identify the \(E_{\alpha}\) with their canonical images in G, and for each \(x\in G\) we shall denote by \(\lambda(x)\) the unique index \(\alpha\in I\) such that \(x\in E_{\alpha}\) . Let \(R\{x,y\}\) denote the following relation between two elements x, y of G: ``there exists an element \(\gamma\in I\) such that \(\gamma\geqslant\alpha=\lambda(x)\) and \(\gamma\geqslant\beta=\lambda(y)\) for which \(f_{\gamma x}(x)=f_{\gamma\beta}(y)\) ''. Then R is an equivalence relation on G. It is clear that R is reflexive and symmetric on G. To show that R is transitive, let \(x\in E_{\alpha}, y\in E_{\beta}, z\in E_{\gamma}\) , and suppose that there exist \(\lambda\in I\) such that \(\lambda\geqslant\alpha,\lambda\geqslant\beta\) and \(f_{\lambda\alpha}(x)=f_{\lambda\beta}(y)\) , and \(\mu\in I\) such that \(\mu\geqslant\beta,\mu\geqslant\gamma\) , and \(f_{\mu\beta}(y)=f_{\mu\gamma}(z)\) . Since I is a directed set, there exists \(v\in I\) such that \(v\geqslant\lambda\) and \(v\geqslant\mu\) ;

by \((\mathbf{LI}_{\mathbf{I}})\) we then have

\[
\begin{array}{r l} f _ {\nu \alpha} (x) & = f _ {\nu \lambda} (f _ {\lambda \alpha} (x)) = f _ {\nu \lambda} (f _ {\lambda \beta} (y)) = f _ {\nu \beta} (y) \\ & = f _ {\nu \mu} (f _ {\mu \beta} (y)) = f _ {\nu \mu} (f _ {\mu \gamma} (z)) = f _ {\nu \gamma} (z), \end{array}
\]

which establishes our assertion. The quotient set \(\mathbf{E} = \mathbf{G} / \mathbf{R}\) is called the direct limit of the family \((\mathbf{E}_{\alpha})_{\alpha \in \mathbf{I}}\) with respect to the family of mappings \((f_{\beta \alpha})\) , and is written \(\mathbf{E} = \lim (\mathbf{E}_{\alpha}, f_{\beta \alpha})\) , or simply \(\mathbf{E} = \lim \mathbf{E}_{\alpha}\) when there is no risk of ambiguity. By abuse of language, the pair \((\overrightarrow{(\mathbf{E}_{\alpha}), (f_{\beta \alpha})})\) (which is usually written \((\mathbf{E}_{\alpha}, f_{\beta \alpha}))\) is called a direct system of sets, relative to the directed set I.

Clearly E is not empty provided at least one of the \(E_{\alpha}\) is not empty. We denote by \(f_{\alpha}\) the restriction to \(E_{\alpha}\) of the canonical mapping f of G onto E = G/R; \(f_{\alpha}\) is called the canonical mapping of \(E_{\alpha}\) into E. If \(\alpha \leqslant \beta\) , we have the relation

\[
f _ {\beta} \circ f _ {\beta \alpha} = f _ {\alpha}\tag{22}
\]

since for each \(x \in \mathbf{E}_{\alpha}\) we have \(f_{\beta \beta}(f_{\beta \alpha}(x)) = f_{\beta \alpha}(x)\) by \((\mathrm{LI}_{\mathbf{I}})\) ; and therefore the elements \(x \in \mathbf{E}_{\alpha}\) and \(f_{\beta \alpha}(x) \in \mathbf{E}_{\beta}\) are congruent mod \(\mathbf{R}\) .

Examples

\begin{enumerate}
\def\labelenumi{(\arabic{enumi})}
\item
  Let A, B be two sets, and let \((\mathrm{V}_{\alpha})_{\alpha\in\mathbf{I}}\) be a family of subsets of A whose index set I is directed, and such that the relation \(\alpha\leqslant\beta\) implies \(V_{\beta}\subset V_{\alpha}\) . Let \(E_{\alpha}\) denote the set of all mappings of \(V_{\alpha}\) into B, and for each pair of indices \((\alpha,\beta)\) such that \(\alpha\leqslant\beta\) let \(f_{\beta\alpha}\) be the mapping of \(E_{\alpha}\) into \(E_{\beta}\) which sends each function \(u\in E_{\alpha}\) to its restriction to \(V_{\beta}\) . It is obvious that the conditions \((\mathrm{LI}_{\mathbf{I}})\) and \((\mathrm{LI}_{\mathbf{II}})\) are satisfied, and the set \(E=\lim E_{\alpha}\) is called the set of germs of mappings of the \(V_{\alpha}\) into B. *The most frequent case is that in which \((\mathrm{V}_{\alpha})\) is the family of neighbourhoods of a subset of a topological space A (General Topology, Chapter I, §6, no. 10). *
\item
  Suppose that, for each \(\alpha \in I\) , \(E_{\alpha}\) is the same set \(F\) and that whenever \(\alpha \leqslant \beta\) , \(f_{\beta \alpha}\) is the identity mapping of \(F\) onto itself. Then there exists a canonical bijection of \(\lim E_{\alpha}\) onto \(F\) . In order to define \(\lim E_{\alpha}\) , we have to form the set \(G\) which is the sum of the family \((E_{\alpha})\) ; \(G\) is therefore the union of a family \((G_{\alpha})\) of mutually disjoint sets, and for each \(\alpha \in I\) there is a canonical bijection \(h_{\alpha}: F \to G_{\alpha}\) . We have next to consider the equivalence relation \(R\) on \(G\) corresponding to the partition \((P_{y})_{y \in F}\) , where \(P_{y}\) is the set of all \(h_{\alpha}(y)\) as \(\alpha\) runs through \(I\) . Clearly \(y \to P_{y}\) is a bijection whose inverse is the bijection required. We shall identify \(F\) with \(\lim E_{\alpha}\) by means of this canonical bijection.
\end{enumerate}

LEMMA 1. Let \((\mathbf{E}_{\alpha}, f_{\beta \alpha})\) be a direct system of sets, \(\mathbf{E} = \lim_{\rightarrow} \mathbf{E}_{\alpha}\) its direct limit, and for each \(\alpha \in \mathbf{I}\) let \(f_{\alpha}: \mathbf{E}_{\alpha} \to \mathbf{E}\) the canonical mapping.

\begin{enumerate}
\def\labelenumi{(\roman{enumi})}
\item
  Let \((x^{(i)})_{1\leqslant i\leqslant n}\) be a finite system of elements of E. Then there exists \(\alpha \in I\) and a finite system \((x_{\alpha}^{(i)})_{1\leqslant i\leqslant n}\) of elements of \(\mathbf{E}_{\alpha}\) such that \(x^{(i)} = f_{\alpha}(x_{\alpha}^{(i)})\) for \(1\leqslant i\leqslant n\) .
\item
  Let \((y_{\alpha}^{(i)})_{1\leqslant i\leqslant n}\) be a finite system of elements of some \(\mathbf{E}_{\alpha}\) . If \(f_{\alpha}(y_{\alpha}^{(i)}) = f_{\alpha}(y_{\alpha}^{(j)})\) for each pair of indices \((i,j)\) , then there exists \(\beta \geqslant \alpha\) such that \(f_{\beta \alpha}(y_{\alpha}^{(i)}) = f_{\beta \alpha}(y_{\alpha}^{(j)})\) for each pair \((i,j)\) .
\item
  By the definition of E there exists for each \(i\) an index \(\beta_{i} \in I\) and an element \(z_{\beta_{i}} \in E_{\beta_{i}}\) such that \(x^{(i)} = f_{\beta_{i}}(z_{\beta_{i}})\) . Take \(\alpha\) such that \(\alpha \geqslant \beta_{i}\) for \(1 \leqslant i \leqslant n\) , and \(x_{\alpha}^{(i)} = f_{\alpha \beta_{i}}(z_{\beta_{i}})\) .
\item
  By the definition of E, for each pair \((i, j)\) there exists \(\gamma_{ij} \in I\) such that \(\gamma_{ij} \geqslant \alpha\) and \(f_{\gamma_{ij}\alpha}(y_{\alpha}^{(i)}) = f_{\gamma_{ij}\alpha}(y_{\alpha}^{(j)})\) . Take \(\beta\) such that \(\beta \geqslant \gamma_{ij}\) for all pairs \((i, j)\) , and use the relations \(f_{\beta\alpha} = f_{\beta\gamma_{ij}} \circ f_{\gamma_{ij}\alpha}\) .
\end{enumerate}

